\section{来源审计与课程主问题}

这讲课并不是在宣布“机器人领域已经有了 GPT 时刻”，而是在提出一份 2023 年的研究配方：能否把机器学习 scaling 的设计原则、不断进步的 Internet-scale language / vision models，以及大规模离线机器人数据组合起来，让真实机器人能力不再靠一项一项手工工程缓慢增长。Ted Xiao 随后用 RT-1、SayCan、Inner Monologue 和 DIAL 展示这份配方如何分别进入 skill learning、planning、feedback 与 data augmentation。

\subsection{本讲如何重建}

课程官网提供了录像与三篇必读论文，却没有公开独立 slide deck。因此本讲以 Stanford Online 1080p 官方录像为母版，对 screen-share 区域每两秒采样一次，从 2,284 帧中筛出 315 个高召回稳定候选，再人工去重为 45 个独立教学状态；最终按核定时间戳从 1920x1080 原视频重新提取全宽帧，避免截断幻灯片右侧。口头动机和问答来自官方 `en-US` 人工字幕；本地旧字幕是 4,183 条重复滚动自动字幕，已由 1,840 条非空人工 caption 替换。

\begin{longtable}{L{0.21\textwidth}L{0.31\textwidth}L{0.38\textwidth}}
\toprule
证据层 & 本地材料 & 在讲义中的职责 \\
\midrule
官方录像 & 1:16:07 视频与 45 张恢复 teaching states & 决定课堂顺序、架构图、结果表、演示与系统地图 \\
官方人工字幕 & `lecture15.en.srt` 与 timed transcript & 恢复 speaker motivation、工程取舍、问答限制与不确定性 \\
课程必读论文 & RT-1、SayCan、Inner Monologue & 核对公式、实验定义、结果边界与模块接口 \\
补充一手论文 & BC-Z 与 DIAL & 解释历史前驱、离线数据和 VLM relabeling \\
审计文件 & manifest、coverage、teacher-voice ledger & 证明视觉与口述节点均有教学去向 \\
\bottomrule
\end{longtable}

\begin{warningbox}{历史边界}
课堂日期是 2023 年 2 月 7 日。RT-2、Open X-Embodiment、PaLM-E、Gemini Robotics 及后来的 vision-language-action scaling 结果不能倒灌为课堂证据。末尾会单列现代 VLA 解读，但与 Ted Xiao 当时的主张分开。
\end{warningbox}

\subsection{先把相近术语分开}

讲者在口头上会把 foundation model、Internet-scale model 和 LLM 宽松混用，但写成讲义时必须区分训练对象、数据来源和系统职责。下面的术语表先建立阅读坐标，后文再给公式和案例。

\begin{longtable}{L{0.22\textwidth}L{0.30\textwidth}L{0.38\textwidth}}
\toprule
术语 & 解决的问题 & 核心机制与课程关系 \\
\midrule
Foundation model & 一个模型能否复用于大量下游任务 & 依赖规模化训练、通用接口和可迁移表示；本讲把它作为目标而非既成事实 \\
Imitation learning & 没有显式 reward 时怎样从示范学 policy & 监督学习 $p(a\mid o,l)$；RT-1 的稳定起点 \\
Offline robot learning & 如何把采集与训练解耦 & 先收集固定数据集，再反复训练多个方法；减少真实机器人在线试错 \\
Tokenization & 如何让图像、语言与动作进入统一序列模型 & 把连续或高维输入压成离散/紧凑 token；RT-1 的互操作接口 \\
Affordance & 当前状态下某项技能是否可执行 & SayCan 用 value / success score 约束 LLM 的语言先验 \\
Closed-loop planning & 执行失败后怎样更新计划 & Inner Monologue 把 scene、clarification 与 success feedback 回灌给 planner \\
Vision-language model & 怎样用 Internet-scale visual semantics 改写 robot data & DIAL 用 VLM 生成更丰富语言标签，而不是直接控制机器人 \\
\bottomrule
\end{longtable}

\subsection{四个贯穿全讲的问题}

读后续系统时应持续追问：机器人数据到底缺数量还是缺多样性；Transformer 的价值来自参数规模还是统一 token interface；LLM 提出的动作怎样被物理可执行性约束；反馈和语言标签又怎样改变下一轮 policy。只有把四个问题放在同一条系统链上，才能看清四篇工作不是互不相干的 demo。

\begin{importantbox}{本讲的系统论命题}
机器人 foundation model 的难点不只是训练一个更大网络，而是让技能数据、实时控制、语言规划、环境反馈和语义标签拥有可组合接口。Language 在这份 2023 recipe 中更像“universal glue”：连接模块，但不能替代每个模块的物理 grounding。
\end{importantbox}

\subsection{本章小结}

本讲使用官方录像、人工字幕和五篇一手论文重建，并把“研究 recipe”与“已证明 foundation model”严格区分。接下来先理解机器人为什么向往 emergence 与 homogenization，再讨论物理世界为何成为最后一块难倒的 domino。

